(1) 实验室动量方程 $dp/dt=F$，初始 $p(0)=0$，故 $p=Ft$。由 $E^2=m^2c^4+p^2c^2$，$v=pc^2/E$ 得 $v(t)=Ft/[m\sqrt{1+(Ft/mc)^2}]$，$\gamma(t)=\sqrt{1+(Ft/mc)^2}$。

(2) 对速度积分：$x(t)=(mc^2/F)[\sqrt{1+(Ft/mc)^2}-1]$。若 $Ft\ll mc$，$x=Ft^2/(2m)+O(t^4)$；若 $t\to\infty$，$v/c\to1$，且 $x=ct-mc^2/F+o(1)$，始终不超光速。

(3) 总能 $E=mc^2\sqrt{1+(Ft/mc)^2}$，动能 $K=E-mc^2$。对时间求导：$\dot E=(F^2tc^2)/E=Fv$，同时 $Fx=mc^2[\sqrt{1+(Ft/mc)^2}-1]=K$。四动量满足 $(E/c)^2-p^2=m^2c^2$，与时间无关；这保证不变静质量。

(4) $d\tau=dt/\gamma$，积分得 $\tau=(mc/F)\operatorname{arsinh}(Ft/mc)$。沿力方向的固有加速度 $\alpha=\gamma^3 dv/dt=\gamma F/m$，随时间增加，不是常量；与经典恒定固有加速度的双曲线世界线不同。$F\to0$ 时 $\tau\to t$，$v\to0$，$x\to0$。
